\documentclass[a4paper,openbib,10pt]{memoir}
\usepackage[utf8]{inputenc}
\usepackage[english]{babel}
\usepackage[T1]{fontenc}
\usepackage{lmodern}
\usepackage{amsmath}
\usepackage{amssymb}
\usepackage{amsthm}
\usepackage{mathtools}
\usepackage{listings}
\usepackage{enumerate}
\usepackage{enumitem}
\usepackage[dvips]{graphicx}
\usepackage{wrapfig}

\usepackage{csquotes}
\usepackage[backend=bibtex,style=alphabetic]{biblatex}
\addbibresource{phdbooks.bib}

% \usepackage[english, status=draft]{fixme}

\usepackage[all,cmtip]{xy}
\entrymodifiers={+!!<0pt,\fontdimen22\textfont2>}

\setlrmarginsandblock{*}{4cm}{*}
\checkandfixthelayout

\hyphenation{to-po-lo-gi-cal}

\theoremstyle{plain}
\newtheorem{theorem}{Theorem}
\newtheorem{lemma}[theorem]{Lemma}
\newtheorem{proposition}[theorem]{Proposition}
\newtheorem{corollary}[theorem]{Corollary}
\newtheorem*{conjecture}{Conjecture}

\theoremstyle{definition}
\newtheorem{definition}[theorem]{Definition}
\newtheorem{example}[theorem]{Example}

\theoremstyle{remark}
\newtheorem{remark}[theorem]{Remark}

\newcommand{\ip}[2]{\left< #1, #2 \right>}
\newcommand{\size}[1]{\left| #1 \right|}

\let\O\undefined

\newcommand{\N}{\mathbb{N}}
\newcommand{\Z}{\mathbb{Z}}
\newcommand{\Q}{\mathbb{Q}}
\newcommand{\R}{\mathbb{R}}
\newcommand{\C}{\mathbb{C}}
\newcommand{\A}{\mathcal{A}}
\newcommand{\delim}{\:}
\newcommand{\X}{X_{m,n}}
\newcommand{\sdel}{\delim\Bigg\vert\delim}
\newcommand{\SUT}[1]{\SU_{#1}/T^{#1-1}}
\newcommand{\GLB}[1]{\GL_{#1}/B_{#1}}
\newcommand{\UT}[1]{\U_{#1}/T^{#1}}
\newcommand{\GL}{\mathrm{GL}}
\newcommand{\U}{\mathrm{U}}
\newcommand{\SU}{\mathrm{SU}}
\newcommand{\SL}{\mathrm{SL}}
\DeclareMathOperator{\O}{O}
\DeclareMathOperator{\SO}{SO}
\DeclareMathOperator{\GS}{GS}
\DeclareMathOperator{\Ext}{Ext}
\DeclareMathOperator{\Hom}{Hom}
\DeclareMathOperator{\im}{im}
\DeclareMathOperator{\Tr}{Tr}
\DeclareMathOperator{\Id}{Id}
\DeclareMathOperator{\supp}{supp}
\DeclareMathOperator{\pr}{pr}
\DeclareMathOperator{\spa}{span}
\DeclareMathOperator{\inv}{inv}
\DeclareMathOperator{\Det}{Det}

\newcommand{\closure}[2][3]{{}\mkern#1mu\overline{\mkern-#1mu#2}}

\newcommand{\set}[1]{\left\{ #1 \right\}}
\newcommand{\com}[1]{\left[#1,#1\right]}

\newcommand{\lie}[1]{\left[ #1 \right]}

\renewcommand{\a}{\alpha}
\renewcommand{\b}{\beta}
\newcommand{\g}{\gamma}
\newcommand{\e}{\varepsilon}
\renewcommand{\d}{\delta}
\renewcommand{\l}{\lambda}

\begin{document}

\frontmatter

\begin{titlingpage}
  % forside centreret i forhold til papiret ikke i forhold til
  % marginerne
  \calccentering{\unitlength}
  \begin{adjustwidth*}{\unitlength}{-\unitlength}
    \begin{adjustwidth}{-1cm}{-1cm} 
      \vspace*{2cm}
      \centering
      \scshape
      
      {\Huge Approximations to the loop space of \newline\newline flag
        manifolds} \\
      
      \vspace*{3\onelineskip}
      
      \includegraphics[width=8cm]{segla1s.ps}    
      
      \vspace*{3\onelineskip}
      
      {\large
        \linespread{1.4}\selectfont
        Progress Report\\
        Simon Stoltze\\
        January 6, 2015\\
        Advisor: Marcel B\"okstedt\\ }
      
      \strut\vfill
      Institut for Matematik \\
      Aarhus Universitet
    \end{adjustwidth}
  \end{adjustwidth*}
  
\end{titlingpage}


%%% Local Variables: 
%%% mode: latex
%%% TeX-master: "main"
%%% End: 


\chapter{Preface}

% Blah blah blah

The following document is part of the midway exam for Ph.D. students
at Aarhus University. It describes the work I have done on my project
so far, and
ends with a small presentation of ideas for future work. Due to the
limited number of pages, the report is meant more as an overview and
will occasionally only sketch a proof or argument instead of giving
all the details.

I would like to thank Thomas Schmidt for proofreading an earlier
version of this report, and my advisor Marcel B\"okstedt for his help
and inspiration during the past two years.

\subsubsection{Conventions and notation}

What follows is a brief description of the notation used in the
report.
For two topological spaces $X$ and $Y$, we will write $X \cong Y$ to
denote that $X$ and $Y$ are homeomorphic and $X\simeq Y$ to denote
that they are homotopy equivalent. Likewise, $f \simeq g$ means that
the maps $f$ and $g$ are homotopic. For subsets, we will use $A
\subset X$ when $A$ is any subset of $X$, and we will generally not
make a distinction between points of $X$ and one-point subsets of
$X$. The set difference of $X$ and $A$ will be written as $X-A$. With
a distinguished basepoint $x_0$ in $X$, the set of based loops in $X$
will be denoted $\Omega(X) = \set{ \g : S^1 \to X \mid \g(1) = x_0}$,
and this comes equipped with the compact-open topology. The one-point
compactification of $X$ is $X^+ = X\cup \set{\infty}$.

We will be working exclusively with topology in the complex numbers,
so the disc $D$ is the unit disc in the complex plane $\C$ and $D_p$
is the unit disc in $\C^p$. The non-zero elements in $\C$ will be
written as $\C^* = \C - 0$. $\GL_m$ will denote the group of
linear automorphisms of $\C^m$ and the standard basis of $\C^m$ is
$e_1,\dots,e_m$. Linear maps will be written as matrices with respect
to the standard basis.

Since we are mainly concerned with algebraic topology, a space will
occasionally be replaced with a homotopy equivalent one without
mention. For example, $T^n$ will denote both the torus $(\C^*)^n
\subset \GL_n$ and $(S^1)^n\subset \C^n$ depending on context, and the
punctured disc $D_p-0$ will sometimes be identified with the sphere
$S^{2p-1} \subset \C^p$. Likewise, we will not distinguish between the
different models of the flag manifold of $\C^m$, given as $\GLB{m}
\cong \UT{m} \cong \SUT{m}$ with $B_m \subset \GL_m$ the upper (or
lower) triangular matrices in $\GL_m$, $T^m \subset \U_m$ the diagonal
matrices in the unitary group and $T^{m-1}\subset \SU_m$ the diagonal
matrices in the special unitary group.

Homology and cohomology groups will always be singular groups
with integer coefficients. If a theorem is mentioned by name without a
reference, it can be found in \cite{hatcher}.

%\vfill % Er allerede s� langt nede som muligt.
\begin{flushright}
  Simon Stoltze \\
  Aarhus University\\
  January, 2015
\end{flushright}



%%% Local Variables: 
%%% mode: latex
%%% TeX-master: "main"
%%% End: 
 % Er nu et preface
\tableofcontents
\chapter{Introduction}

% Noget om starten, hvorfor dette er interessant. Se bagside af blok
% Noget om notation... eg \subset, D-0, hvad er D_p
% Alting komplekst, s� \GL_m giver mening etc. M�ske ogs� de
% forskellige versioner af flag-mangfoldigheder.
% Og noget med at T=\C^*....

This document gives a description of a family of spaces $\X\subset
\C^{mn}$, indexed by two natural numbers $m$ and $n$, and a quotient
of these spaces, $Y_{m,n} = \X / T^n$. The spaces were
originally defined as the zero-set of a section on a vector bundle of
a Grassmann manifold, and were studied to understand the properties of
families of Calabi-Yau manifolds. However, the spaces turned out to
have other properties that made them interesting in their own
right. The most important property is that they can be used to form a
space that is homotopy equivalent to a space that appears in physics,
namely the loop space $\Omega(\SUT{m})$. This hints at a way of
extracting information about the loop space by working with the
simpler spaces $Y_{m,n}$.

The spaces have been studied by Marcel B\"okstedt, and some of his
results are gathered in \cite{bokstedt}. In particular it includes the
following theorem that covers the special case where $m=2$:
\begin{theorem}
  The homology of $Y_{2,n}$ is
  \[ H_q(Y_{2,n}) =
  \begin{cases}
    \Z & 0 \leq q \leq n \\
    0 & \text{otherwise}
  \end{cases} \]
  There is a stabilization map $s : Y_{2,n} \to Y_{2,n+1}$ that
  induces an isomorphism
  \[ s_* : H_q(Y_{2,n}) \to H_q(Y_{2,n+1}) \]
  for $q \leq n$. There is also a map $f : Y_{2,n} \to
  \Omega(\mathrm{SU_2}/T^1)$ that induces an isomorphism in homology in
  degree $q\leq n$.
\end{theorem}

The aim of the following is to try and extend these results to the
general case of $Y_{m,n}$, with $m\geq 2$.

This progress report is divided into three main parts.
The first chapter gives an introduction to the spaces we are
interested in and lists some of the properties that will be used
in the other chapters.
The second chapter is mainly concerned with an explicit calculation of
the cohomology of the spaces $X_{3,3}$ and $Y_{3,3}$. This is done by
applying the spectral sequence coming from a filtration of $X_{3,3}$,
computing the various relative cohomology groups needed, and
finding the differentials of the sequence.
The third and final chapter concerns the relation to the loop space
described above. Here it is shown that by taking a limit of the spaces
$Y_{m,n}$, we get a space homotopy equivalent to the loop space. The
chapter concludes with a brief overview of plans for future studies.

%%% Local Variables: 
%%% mode: latex
%%% TeX-master: "main"
%%% End: 
 % Snak med Marcel. Derudover, tag fra
                       % projektbeskrivelsen.
\mainmatter


\chapter{The space $\X$}

% Snak snak. M�ske deles op i to?
% Noget om starten, hvorfor dette er interessant. Se bagside af blok
% Ogs� noget om rummene, definition, symmetrier (specielt X_{m,n}
% \cong X_{n,m}, m�ske filtreringen?
This chapter will introduce the spaces to be studied, along with a
brief listing of various useful structure.

\begin{definition}
  For natural numbers $m$ and $n$, the space $X_{m,n} \subset
  \C^{mn}$ is defined as
  \[ X_{m,n} = \set{(a_1,\dots,a_n) \in (\C^m)^n \delim\Bigg\vert\delim
    \begin{matrix}
      \text{Any } m \text{ subsequent vectors in } \\
      (e_1,\dots,e_m,a_1,\dots,a_n,e_1,\dots,e_m) \\
      \text{ are linearly independent.}
    \end{matrix} } \]
  For any two elements $X = [x_1,\dots,x_m]$ and $Y = [y_1,\dots,y_m]$
  in $\GL_m$, define the space $X_{m,n}(X,Y)$ as
  \[ X_{m,n}(X,Y) = \set{(a_1,\dots,a_n) \in (\C^m)^n
    \delim\Bigg\vert\delim
    \begin{matrix}
      \text{Any } m \text{ subsequent vectors in } \\
      (x_1,\dots,x_m,a_1,\dots,a_n,y_1,\dots,y_m) \\
      \text{ are linearly independent.}
    \end{matrix} } \]
  The special case $X=\Id$ will be denoted $\X(Y)$.
  Elements $(a_1,\dots,a_n)\in \X$ will be identified
  with the $m\times n$ matrix $A = [a_1,\dots,a_n]$ without mention.
\end{definition}

Since the space $\X(X,Y)$ only depends on linear independence of
subsequent vectors, it can be described by giving the two flags in
$\C^m$ defined from the columns of the matrices
$X=\left[x_1,\dots,x_m\right]$ and $Y=\left[y_1,\dots,y_m\right]$:
\begin{align*}
  \mathrm{Fl_R}(X) &= \Big(\spa(x_m) \subset \spa(x_{m-1},x_m) \subset
    \dots \subset \C^m\Big) \\
  \mathrm{Fl_L}(Y) &= \Big(\spa(y_1) \subset \spa(y_1,y_2) \subset
    \dots \subset \C^m\Big)
\end{align*}
Multiplying with $X^{-1}$ on each vector in $\X(X,Y)$ defines a
homeomorphism 
\[ \X(X,Y)\cong \X(\Id,X^{-1}Y) = \X(X^{-1}Y) \]
Hence we only need to consider these spaces.

Note that we can also define the space $X_{m,n}$ as the preimage of
$(\C^*)^{m+n-1}$ under the map $\Det : (\C^m)^n \to \C^{m+n-1}$,
given by
\begin{align*}
  \Det(a_1,\dots,a_n) = \Big( &\det(e_2,\dots,e_m,a_1),\\
  &\det(e_3,\dots,e_m,a_1,a_2),\\
  &\dots, \\
  &\det(a_n,e_1,\dots,e_{m-1}) \Big) 
\end{align*}
where $\det$ is the determinant map. There is a similar description
of $X_{m,n}(Y)$ as the preimage under a map $\Det_{Y}$. Since the map
$\Det_Y$ is continuous, this shows that the space $\X(Y)$ is open in
$\C^{mn}$.

\section{Stabilization}

We would like to relate the spaces $X_{m,n}$ and $X_{m,n+1}$. To do
this, consider a lower-triangular $m\times m$ invertible matrix
$L$. This preserves the right flag $\mathrm{Fl_R}(\Id)$, so it defines
a homeomorphism $\tilde L : \X(Y) \to \X(LY)$ by
\begin{align*}
  \tilde L(a_1,\dots,a_n) &= (La_1,\dots,La_n)
\end{align*}

By using the Bruhat decompositon of the general linear group,
any matrix $Y$ can be written uniquely as a product $L \sigma U$,
where $L$ is an invertible lower triangular matrix, $U$ is an
invertible upper triangular matrix, and $\sigma$ is a
permutation matrix, see \cite[Example~1.2.11]{bjorner} or
\cite[Proposition~4.5]{hiller}. Since the flag $\mathrm{Fl_L}(\sigma
U)$ is preserved by right multiplication with upper triangular
matrices, this only depends on the permutation $\sigma$ and not on
$U$. All of this together shows the following lemma:

\begin{lemma}
  For $Y\in \GL_m$ with $Y = L\sigma U$, the space $\X(Y)$ is
  homeomorphic to $\X(\sigma)$.
\end{lemma}

Using this lemma, we only need to understand the spaces
$\X(\sigma)$ to understand all of the spaces $\X(Y)$. More importantly,
it turns out that the space $X_{m,n+1}(\sigma)$ is given as a union of
spaces homeomorphic to a torus times $\X(\tau)$, for various
permutations $\tau$. This is summarized in the following lemma.

\begin{lemma}
  Any choice of indices $I= (i_1<\dots<i_k) \subset \set{1,\dots,m}$
  gives a subspace of $X_{m,n+1}(\sigma)$:
  \[ X_{m,n+1}^I(\sigma) = \set{(a_1,\dots,a_{n+1}) \in
    X_{m,n+1}(\sigma) \sdel 
    \begin{matrix} 
      (a_{n+1})_{i_j} \neq 0\; \forall i_j \in I, \\
      (a_{n+1})_j = 0 \;\forall j \not\in I
    \end{matrix} } \]
  This space is empty if $I$ does not contain the number $\sigma(m)$,
  so we will assume that $\sigma(m) \in I$ but otherwise leave it out
  of the notation.
  
  These subspaces cover $X_{m,n+1}(\sigma)$. There is a map $\varphi$,
  defined for an indexing set $I$ and a permutation $\sigma$, which
  gives a new permutation such that
  \[ X_{m,n+1}^I(\sigma) \cong \X(\varphi(I,\sigma)) \times
  (\C^*)^{\size{I}} \]
  If we define the permutation $\widehat \sigma = \sigma \cdot
  \big(m\; m-1\; \dots \; 1\big)$ and the indexing set is $I = (i_1 <
  \dots < i_k)$, then $\varphi$ is given by
  \[ \varphi\big(I,\sigma\big) = \left(\prod_{j \in J} \big(
      i_j \; \sigma(m) \big) \right) \widehat \sigma \]
  where the product is over the set
  \[ J = \set{ j \in\set{1,\dots,k} \delim\mid\delim i_j <
    \sigma(m),\; \widehat\sigma^{-1}(i_j) > \widehat\sigma^{-1}(i_r) \,
    \forall r < j } \]
\end{lemma}

The proof of this lemma relies on taking $X_{m,n+1}^I(\sigma)$ and
multiplying with a lower triangular matrix to reduce it to something
in $\X(\tau)$. By considering the matrix used, the formula
for $\tau = \varphi(I,\sigma)$ given in the lemma can be found. In
practice, it is
often easier to do the reduction by hand rather than using the lemma,
but the formula can be useful. For example, it is clear that
$\varphi(I,\sigma)$ is given as a product of transpositions
and $\varphi\big((\sigma(m)), \sigma\big)$. By
considering exactly which transpositions, we can show that for all $I$
we have
\[ \varphi(I,\sigma) \leq \varphi\big((\sigma(m)), \sigma\big) \] 
in the Bruhat ordering on the symmetric group. The Bruhat
order is defined in \cite{bjorner}. It also follows that for all
$\sigma$, 
\[ \varphi\big((1<2<\dots<m),\sigma\big) = \Id \]
With this we can identify $X_{m,n}$ with the subspace
of $X_{m,n+1}$ where the last column has a 1 on each entry:
\[ X_{m,n} \cong \set{ (a_1,\dots,a_{n+1}) \in X_{m,n+1} \mid
  a_{n+1} = (1,\dots,1) } \]
The identification is given by the map $s : \X \to X_{m,n+1}$,
\[ s(a_1,\dots,a_n) = (La_1,\dots,La_n,Le_1) \]
where $L$ is the lower triangular matrix with all entries on or below
the diagonal equal to one:
\[ L =
\begin{pmatrix}
  1 & 0 & \dots & 0 \\
  \vdots & \ddots & \ddots & \vdots \\
  \vdots &  & \ddots & 0 \\
  1 & \dots & \dots & 1
\end{pmatrix} \]

These maps give a directed system of spaces,
\[ \xymatrix{ X_{m,1} \ar[r]^s & X_{m,2} \ar[r]^s & \dots \ar[r]^s &
  X_{m,n} \ar[r]^s & X_{m,n+1} \ar[r]^s & \dots } \]
and we can take the direct limit of this system,
\[ X_{m,\infty} = \varinjlim_n \X \]
The direct limit is the disjoint union of all the spaces $\X$, with the
identification $A \sim B$ if there are $j,k$ such that
$s^k(A) = s^j(B)$. The topology is the finest such that the maps $\X
\to X_{m,\infty}$ are all continuous.

\section{Symmetries}

% Skriv noget her og byt rundt p� r�kkef�lge?
% Noget om virkninger af \C^* p� \X, e.g. virk p� s�jler, r�kker,...
There are various symmetries and group actions worth keeping in mind
when working with these spaces. Some of these will be described here.

Since we are interested in linear independence of vectors in $\C^m$,
the group of complex units $\C^*$ acts on $\X(\sigma)$ in various
ways. For any index $i$, we can define a group action $l_i$ by scaling
the $i$\textsuperscript{th} vector:
\begin{align*}
  l_i : \C^* \times \X(\sigma) &\to \X(\sigma) \\
  \big(\lambda , (a_1,\dots,a_n)\big) &\mapsto
  \big(a_1,\dots,a_{i-1},\lambda a_i,a_{i+1},\dots,a_n\big) 
\end{align*}
These actions commute with each other, so we get an action $l$ of
$(\C^*)^n$ on $\X(\sigma)$, given by
\[ l\big((\l_1,\dots,\l_n),(a_1,\dots,a_n)\big) =
\big(\l_1a_1,\dots,\l_na_n\big) \]

Likewise, we can also multiply the rows of $A\in
\X(\sigma)$ by non-zero complex numbers. This defines an action $l'$
of $(\C^*)^m$ on $\X(\sigma)$. The two actions together satisfy the
relation 
\[ l\Big((\l,\dots,\l),l'\big((\l^{-1},\dots,\l^{-1}), A\big)\Big) =
A \]
since multiplying all rows by $\l^{-1}$ while multiplying all columns
by $\l$ is the same as doing nothing.

For the special case of $\X = \X(\Id)$, there is an additional
symmetry that can be useful to know. By thinking of $A \in \X$ as a
matrix, it is possible to transpose it. Checking the various
determinants used in the definition of $\X$ show that they do not
change after transposition. Hence transposing defines a homeomorphism
$T : X_{m,n} \to X_{n,m}$, allowing us to switch the order of $m$
and $n$.

\section{The quotient space}

% Rummene Y_{m,n}. 
While the above definition of $\X$ gives an open subset of $\C^{mn}$,
it has some issues that makes the space hard to
work with. As we shall see shortly, the fundamental group and the
cohomology ring of $\X$ both become larger as $n$ grows.
This complicates things slightly, but it can be fixed by taking the
quotient of a group action.
\begin{definition}
  The space $Y_{m,n}$ is the quotient space
  \[ Y_{m,n} = \X / T^n \]
  where $T^n$ acts on $X_{m,n}$ by scaling
  the columns, as defined above:
  \[ (\lambda_1,\dots,\lambda_n)\cdot (a_1,\dots,a_n) =
  (\l_1a_1,\dots,\l_na_n) \]
\end{definition}

These spaces retain some of the structure described above. For
example, the stabilization map $s$ is linear in $A$, which means it
respects the group action.
\[ s\big((\l_1,\dots,\l_n) \cdot A\big) = (\l_1,\dots,\l_n,1)\cdot
s(A) \] 
Hence $s$ descends to a stabilization map $s : Y_{m,n} \to
Y_{m,n+1}$ and there is a direct limit,
\[ Y_{m,\infty} = \varinjlim_{n} Y_{m,n} \]
which can be thought of as $X_{m,\infty}/T^\infty$.

The space $Y_{m,n}$ is homeomorphic to the subspace of $\X$ consisting
of elements with the last $n$ coordinates of the determinant map
$\Det$ equal to 1:
\[ Y_{m,n} \cong \set{ A \in \X \delim\big\vert\delim \Det(A) \in
  (\C^*)^{m-1}\times \set{1}^n }\]
The group action gives an identification,
\[ \X \cong Y_{m,n}\times (\C^*)^n \]
This shows the claim about the fundamental group of $\X$, but it can
also be used to get information about $Y_{m,n}$ from $\X$ or
vice versa. For example, it allows us to calculate the cohomology of
one when we know the cohomology of the other, by using the
K\"unneth formula to add or remove the cohomology of $(\C^*)^n$.



%%% Local Variables: 
%%% mode: latex
%%% TeX-master: "main"
%%% End: 
 %  Introducer rummene her. Husk stabilisering, symmetri,
              %  noget med reduktion X_{m,n,\set{...}} \cong
              %  X_{m,n-1}\times \C^*^r
              %  Og noget med kun afhængig af permutationen


\chapter{Cohomology calculations for $X_{3,3}$}

In this section we will calculate the cohomology of the space
$X_{3,3}$, which is the set:
\[ X_{3,3} = \set{
  \begin{pmatrix}
    a & b & g \\
    c & d & h \\
    e & f & i
  \end{pmatrix} \in \C^9 \sdel 
  \begin{matrix}
    a\neq 0,\, ad-bc \neq 0, \\
    adi+beh+cfg-deg-bci-afh \neq 0, \\
    di-fh \neq 0,\, i \neq 0
  \end{matrix} } \]

The way we will calculate this is by using the spectral sequence for
a filtration of $X_{3,3}$. Some of the calculations have
been done in Sage and the relevant source code can be found at
\url{http://home.math.au.dk/sstoltze/del-a.html}.
% Some description of this space, properties, something?

\section{A filtration on $\X$}

Before we start the calculation, we need some general setup. To work
with the space $X_{m,n}$ and compute the cohomology, it will
be beneficial to build it up from simpler subspaces. This is done by
considering a filtration.

For $0 \leq p \leq m-1$, let $F_p$ be the subset
\[ F_p=\set{ (a_1,\dots,a_n) \in X_{m,n} \mid a_n \text{ has at most } p
  \text{ zeroes}.} \]
$F_p$ is open in $\X$ since we stay inside the subspace when we
vary the coordinates of a point $A\in F_p$, as long as we do not
introduce new zeroes in the last column.
Together, the subspaces form an increasing sequence:
\[ F_0 \subset F_1 \subset \dots \subset F_{m-1} = X_{m,n} \]

The general method of computing the homology and cohomology of
$X_{m,n}$ is by using the spectral sequence of this
filtration, as described in \cite{hatcherss} and
\cite{mccleary}. Here we will focus on
cohomology, but similar arguments could of course be applied to
homology.

The spectral sequence has $E_1$-page given by
\[ E_1^{p,q} = H^{p+q}(F_p,F_{p-1}) \]
with differentials given by the boundary map in the long exact
sequence of the triple $(F_{p+1},F_p,F_{p-1})$:
\[ d_1 : H^{p+q}(F_p,F_{p-1}) = E_1^{p,q} \to E_1^{p+1,q} =
H^{p+q+1}(F_{p+1},F_p) \]

To use the spectral sequence, we need to compute the cohomology of the
pair $(F_p,F_{p-1})$. To calculate this we will use excision, so we
would like to find an open subset of $F_p$ that contains the
complement of $F_{p-1}$. For any choice of exactly $p$ 
different indices $\set{v_1,\dots,v_p}\subset \set{1,\dots,m-1}$,
define $X_{\set{v_1,\dots,v_p}}$ to be the points $(a_1,\dots,a_n)\in
F_p$ which has $(a_n)_{v_i} = 0$ for all $i$. Since $\X$ is open, we
can get an open set in $F_p$ by inserting a small disc on these
zero-entries. More precisely, the set we get is
\[ X_{\set{v_1,\dots,v_p}} \times D_p \cong \set{ A\in F_p \sdel 
  \begin{matrix} \displaystyle 
    \size{((a_n)_{v_1},\dots,(a_n)_{v_p})} <
    \min_{i\not\in \set{v_1,\dots,v_p}} \size{(a_n)_i}, \\
    \left(a_1,\dots,a_{n-1},a_n-t\sum_{i} (a_n)_{v_i}e_{v_i}\right)\in
    F_p \forall t \in [0,1]
  \end{matrix} } \]
This is an open set in $F_p$, with $X_{\emptyset} = F_0$, and it
contains all points $(a_1,\dots,a_n)\in F_p$ where the entries of
$a_n$ indicated by $\set{v_1,\dots,v_p}$ are zero. For a different
choice of indices, the condition that the radius of the disc is
smaller than the norm of the non-zero entries in $a_n$ ensures
disjointness, so the intersection is empty:
\[ \left(X_{\set{v_1,\dots,v_p}}\times D_p\right) \cap \left(
X_{\set{u_1,\dots,u_p}}\times D_p\right) = \emptyset \quad \text{if }
\set{v_1,\dots,v_p} \neq \set{u_1,\dots,u_p} \]

The intersection of $\X\times D_p$ and $F_{p-1}$ corresponds to
removing 0 from the disc, since one of the entries inserted on the
last column has to be non-zero.
\[ \left(X_{\set{v_1,\dots,v_p}}\times D_p \right) \cap F_{p-1} =
X_{\set{v_1,\dots,v_p}} \times \left(D_p-0\right) \]
This means that the disjoint union of all of these form an open set in
$F_p$ with
\[ F_p = F_{p-1}\cup \left(\coprod_{\set{v_1,\dots,v_p}}
X_{\set{v_1,\dots,v_p}}\times D_p\right) \]
and by excision, the inclusion of pairs
\[ \xymatrix{ \displaystyle \coprod_{\set{v_1,\dots,v_p}}
  X_{\set{v_1,\dots,v_p}}\times (D_p,D_p-0) \ar@{^(->}[r] &
  (F_p,F_{p-1}) } \]
induces an isomorphism of cohomology groups,
\begin{align*}
  E_1^{p,q} &= H^{p+q}(F_p,F_{p-1}) \\
  &\cong H^{p+q}\left(\coprod_{\set{v_1,\dots,v_p}}
    X_{\set{v_1,\dots,v_p}}\times (D_p,D_p-0)\right) \\   &\cong
  \bigoplus_{\set{v_1,\dots,v_p}}
  H^{p+q}(X_{\set{v_1,\dots,v_p}}\times (D_p,D_p-0)) \\   &\cong
  \bigoplus_{\set{v_1,\dots,v_p}}
  H^{q-p}(X_{\set{v_1,\dots,v_p}})\otimes H^{2p}(D_p,D_p-0) 
\end{align*}
The last isomorphism is given by the K\"unneth formula, and the fact
that the pair $(D_p,D_p-0)$ only has non-zero cohomology in degree
$2p$:
\[ H^q(D_p,D_p-0) \cong \widetilde H^{q-1}(D_p-0) = \widetilde
H^{q-1}(S^{2p-1}) =
\begin{cases}
  \Z & q = 2p \\
  0 & \text{otherwise}
\end{cases} \]


\section{The spectral sequence of $X_{3,3}$}

We can apply the above filtration to get a spectral sequence of
$X=X_{3,3}$. The groups we need to compute for the first page are:
\[ E_1^{p,q} \cong
\begin{cases}
  H^{q}(X_\emptyset) & p=0 \\
  H^{q+1}(X_{\set{1}} \times (D,D-0)) \oplus H^{q+1}(X_{\set{2}} \times
  (D,D-0)) & p=1 \\
  H^{q+2}(X_{\set{1,2}}\times (D_2,D_2-0)) & p = 2 \\
  0 & \text{otherwise}
\end{cases} \]
The four spaces that appear will be handled separately.

\subsection{The cohomology of $X_{\set{1,2}}$}

Consider the space $X_{\set{1,2}}$. This consists of the elements of
$X_{3,3}$ that have two zeroes in the last column:
\[ X_{\set{1,2}} = \set{
  \begin{pmatrix}
    a & b & 0 \\
    c & d & 0 \\
    e & f & i
  \end{pmatrix} \sdel
  \begin{matrix}
    a \neq 0, ad-bc \neq 0 \\
    adi-bci \neq 0 \\
    di \neq 0, i \neq 0
  \end{matrix}
} \]

The space is homeomorphic to $(\C^*)^3\times \C^2
\times \set{(b,c)\in\C^2 \delim \vert \delim bc\neq 1}$, with the
homeomorphism given by
\begin{align*}
  X_{\set{1,2}} &\to (\C^*)^3\times \C^2
  \times \set{(b,c)\in\C^2 \delim \vert \delim bc\neq 1} \\
  \begin{pmatrix}
    a & b & 0 \\
    c & d & 0 \\
    e & f & i
  \end{pmatrix} &\mapsto \left(a,d,i,\frac{e}{i},\frac{f}{i},
    \frac{b}{a}, \frac{c}{d}\right)
\end{align*}

The space $\set{(b,c)\delim \vert\delim bc\neq 1}$ will be denoted
$Y$. To get the cohomology of this space, consider it as a subspace of
the one-point compactification $S^4 = \C^2\cup\set{\infty}$. Then the
complement is:
\[ S^4-Y = \set{(b,c)\in \C^2 \mid bc=1}\cup\set{\infty} \]
This can be identified with the wedge sum $S^1\vee S^2$ by considering
the subspace
\[ \C^* \cong \set{\left(b,\frac{1}{b} \right) \sdel b\in \C^*} =
(S^4-Y)-\set{\infty} \]
The point at infinity is glued to this space when either $b$ or $c$
is large, which for our copy of $\C^*$
corresponds to $b$ large or $b$ near 0. Hence what we get is
$S^2=\C\cup\set{\infty}$ with the identification $0 \sim \infty$. This
is homotopy equivalent to $S^1\vee S^2$.

Now we can apply Alexander duality to this space to compute
\[ \widetilde H_q(Y) \cong \widetilde H^{3-q}(S^4-Y) =
\begin{cases}
  \Z & q \in \set{1,2} \\
  0 & \text{otherwise}
\end{cases} \]
and hence
\[ H^q(Y) \cong
\begin{cases}
  \Z & q \in \set{0,1,2} \\
  0 & \text{otherwise}
\end{cases} \]

From all of this, we can calculate the cohomology of $X_{\set{1,2}}$
\[ H^q(X_{\set{1,2}}) =
\begin{cases}
  \Z & q = 0 \\
  \Z^4 & q = 1 \\
  \Z^7 & q = 2 \\
  \Z^7 & q = 3 \\
  \Z^4 & q = 4 \\
  \Z & q = 5 \\
  0 & \text{otherwise}
\end{cases} \]

Since we know the product structure from the K\"unneth formula, we can
name the generators and get the ring structure,
\[ H^*(X_{\set{1,2}}) = \bigwedge\nolimits^*[a_{00},d_{00},i_{00}]\otimes
H^*(Y) \]
The two generators of $Y$ will be denoted $y_{00,1}$ for the degree
one generator and $y_{00,2}$ for the degree two generator. All other
generators have degree one. The ring
structure is the usual one on a tensor product, with the added
relation $y_{00,1}\cup y_{00,2} = 0$. The element $y_{00,1}\in H^1(Y)$
is given by the map
\[ Y\ni (b,c) \mapsto 1-bc\in \C^* \]
in cohomology, since there is another map
\[ \C^* \ni \l \mapsto (1,1-\l) \in Y \]
and the composition is the identity on $\C^*$. Applying $H^1$ shows
that the two maps give isomorphisms in degree 1.


\subsection{The cohomology of $X_{\set{1}}$}

In $X_{\set{1}}$ we have the first entry of the last column equal to
zero, which is the space
\[ X_{\set{1}} = \set{
  \begin{pmatrix}
    a & b & 0 \\
    c & d & h \\
    e & f & i
  \end{pmatrix} \sdel
  \begin{matrix}
    a \neq 0, ad-bc \neq 0 \\
    adi + beh - bci - afh \neq 0 \\
    di -fh \neq 0, h \neq 0, i \neq 0 
  \end{matrix}
} \]

This space is homeomorphic to $(\C^*)^5 \times \C\times Y$:
\[ \begin{pmatrix}
  a & b & 0 \\
  c & d & h \\
  e & f & i
\end{pmatrix} \mapsto \left(
a, i, h, \frac{d}{h}-\frac{f}{i}, \frac{b}{h}-\frac{bc}{ah},
\frac{e}{i}, \frac{bhi}{adi-afh},\frac{c}{h}-\frac{e}{i}
\right) \] 
Since the cohomology of $Y$ is already known, applying the K\"unneth
formula gives
\[ H^q(X_{\set{1}}) =
\begin{cases}
  \Z & q = 0 \\
  \Z^6 & q = 1 \\
  \Z^{16} & q = 2 \\
  \Z^{25} & q = 3 \\
  \Z^{25} & q = 4 \\
  \Z^{16} & q = 5 \\
  \Z^6 & q = 6 \\
  \Z & q = 7 \\
  0 & \text{otherwise}
\end{cases} \] 
As before, the ring structure is given by
\[ H^*(X_{\set{1}}) =
\bigwedge\nolimits^*[a_{01},d_{01},f_{01},h_{01},i_{01}] \otimes H^*(Y)\]

\subsection{The cohomology of $X_{\set{2}}$}

The space is
\[ X_{\set{2}} = \set{
  \begin{pmatrix}
    a & b & g \\
    c & d & 0 \\
    e & f & i
  \end{pmatrix} \sdel
  \begin{matrix}
    a \neq 0, ad-bc \neq 0 \\
    adi +cfg -bci - deg \neq 0 \\
    di \neq 0, g\neq 0, i \neq 0 \end{matrix}
} \]
which can also be identified with $(\C^*)^5\times \C\times Y$ using
the homeomorphism
\[ \begin{pmatrix}
  a & b & g \\
  c & d & 0 \\
  e & f & i
\end{pmatrix} \mapsto \left(\frac{a}{g},d,i,g, 1+\frac{cfg}{adi}-
  \frac{eg}{ai}-\frac{bc}{ad}, \frac{f}{i},\frac{b}{a}, \frac{c}{d}
\right)\]
Since we already know the cohomology of this space, we know
\begin{align*}
  H^q(X_{\set{2}}) = 
  \left(\bigwedge\nolimits^*[a_{10},d_{10},f_{10},g_{10},i_{10}]\otimes
  H^*(Y)\right)^q &=
  \begin{cases}
    \Z & q = 0 \\
    \Z^6 & q = 1 \\
    \Z^{16} & q = 2 \\
    \Z^{25} & q = 3 \\
    \Z^{25} & q = 4 \\
    \Z^{16} & q = 5 \\
    \Z^6 & q = 6 \\
    \Z & q = 7 \\
    0 & \text{otherwise}
  \end{cases}
\end{align*}

\subsection{The cohomology of $X_{\emptyset}$}

This space is slightly more complicated than the others. The space
consists of the elements in $X_{3,3}$ where all entries of the last
vector are non-zero.
\[ X_{\emptyset} = \set{
  \begin{pmatrix}
    a & b & g \\
    c & d & h \\
    e & f & i
  \end{pmatrix} \sdel
  \begin{matrix}
    a \neq 0, ad-bc \neq 0, di-fh \neq 0, \\
    adi+beh+cfg-deg-bci-afh \neq 0, \\
    g\neq 0, h\neq 0, i\neq 0
  \end{matrix} }\]
The space is homeomorphic to $(\C^*)^7\times \set{ (x,y,z)\in \C^3
  \mid y-z-xyz = 1}$, with the homeomorphism given by
\begin{align*}
  \begin{pmatrix}
    a & b & g \\
    c & d & h \\
    e & f & i
  \end{pmatrix} \mapsto \bigg(&g, h, i, \frac{a}{g},
    \frac{f}{i}-\frac{d}{h},\\
  &D=\frac{ad-bc}{ah}, \\
  &C=\frac{c}{h}-\frac{a}{g}+\frac{cdi-deh}{fh^2-dhi} +
  \frac{beh-bci}{fgh-dgi},\\
  &\left.\frac{bC}{aD},\frac{c}{hC}-\frac{a}{gC},
    \frac{D}{C}\frac{eh-ci}{fh-di}\right)
\end{align*}

To get the cohomology, we need to compute the cohomology of the space:
\[ Z=\set{(x,y,z)\in\C^3 \mid y-z-xyz = 1} \]
This is the same as the cohomology of $S^2$, but this is hard to see
directly. Instead consider the space
\[ \widetilde Z = \set{(b,c,d,e) \in \C^4 \mid d\neq
  bc, c\neq de} \]
This space is homeomorphic to $Z\times (\C^*)^2$:
\begin{align*}
  (b,c,d,e) \mapsto \left(
    \frac{b(c-de)}{d-bc}, \frac{c}{c-de}, \frac{e(d-bc)}{c-de},
  d-bc, c-de\right)
\end{align*}
We will use Alexander duality to compute the cohomology as we did for
$Y$. Consider the complement in $S^8$,
\[ S^8-\widetilde Z = \set{(b,c,d,e) \mid d = bc} \cup \set{(b,c,d,e)
  \mid c = de} \cup \set{\infty} \]
and denote the two sets as $W_1$ and $W_2$. Then
\begin{align*}
  W_1\cap W_2 &= \set{(b,c,d,e) \mid d=bc,c=de}\cup\set{\infty}
\end{align*}
We glue the point at infinity to this space when $c$ is large, but
this is the same as either $d$ or $e$ being large. Hence there is a
homeomorphism,
\begin{align*}
  W_1\cap W_2 &\cong \set{(b,d,e) \mid d=bde}\cup\set{\infty} \\
  &= \set{(b,d,e)\mid 0 = (be-1)d}\cup\set{\infty} \\
  &= \set{(b,d,e) \mid d = 0} \cup \set{(b,d,e) \mid
    be=1}\cup\set{\infty}
\end{align*}

These new spaces have computable homology, since they can be
identified with spaces built from spheres. The first is relatively
easy,
\[ Z_1 = \set{(b,d,e)\mid d = 0}\cup \set{\infty} =
\set{(b,0,0,d)}\cup\set{\infty} \cong S^4 \]
The second is slightly harder.
\[ Z_2 =\set{(b,d,e)\mid be=1}\cup\set{\infty}=
\set{\left(b,\frac{d}{b},d,\frac{1}{b} \right) }\cup\set{\infty} \]
This set is homeomorphic to the one-point compactification of
$\set{(b,d)\in \C^*\times \C}$. By \cite[Chapter~2.3]{dupontk}, this
is given by the smash product of the one-point compactifications,
\[ Z_2\simeq (\C^*\times\C)^+ \cong (\C^*)^+\wedge\C^+ \cong (S^1\vee
S^2) \wedge S^2 = S^3 \vee S^4 \]
The intersection of these spaces is
\[ Z_1\cap Z_2 = \set{\left(b,0,0,\frac{1}{b}\right)}\cup\set{\infty}
\simeq S^1\vee S^2 \]

From all of this, we get a commutative diagram of inclusions,
\[ \xymatrix{
  S^1\vee S^2 \ar[r]\ar[d] & S^3\vee S^4 \ar[d] \\
  S^4 \ar[r] & W_1\cap W_2
} \]
The two maps from $S^1\vee S^2$ are zero in cohomology (when the
degree is greater than zero) since in each
degree either the domain or the codomain is zero. By applying
Mayer-Vietoris, we get
\[ H^q(W_1\cap W_2) =
\begin{cases}
  \Z & q=0 \\
  0 & q = 1 \\
  \Z & q = 2 \\
  \Z^2 & q = 3 \\
  \Z^2 & q =4 \\
  0 & \text{otherwise}
\end{cases} \]
where the case $q=0$ follows from $W_1\cap W_2$ being connected
and $q = 1$ follows from writing out the relevant part of the
long exact sequence,
\[ \xymatrix{ H^0(S^3\vee S^4)\oplus H^0(S^4) \ar@{->>}[r] &
H^0(S^1\vee S^4) \ar[r] & H^1(W_1\cap W_2) \ar[r] & 0} \]

$W_1$ and $W_2$ are both a copy of $S^6$ since we can get one of
the coordinates from two of the others. Hence we get another diagram,
\[ \xymatrix{
  W_1\cap W_2 \ar[r]\ar[d] & S^6 \ar[d] \\
  S^6 \ar[r] & S^8-\widetilde Z
} \]
and again the map from the intersection is zero. Repeating the above,
we can compute:
\[ H^q(S^8-\widetilde Z) =
\begin{cases}
  \Z & q = 0 \\
  0 & q = 1 \\
  0 & q = 2 \\
  \Z & q = 3 \\
  \Z^2 & q = 4 \\
  \Z^2 & q = 5 \\
  \Z^2 & q =6 \\
  0 & \text{otherwise}
\end{cases} \]
Applying Alexander duality then gives
\[ H^q(\widetilde Z) =
\begin{cases}
  \Z & q = 0 \\
  \Z^2 & q = 1 \\
  \Z^2 & q = 2 \\
  \Z^2 & q = 3 \\
  \Z & q = 4 \\
  0 & \text{otherwise}
\end{cases} \]
from which we get that the cohomology of $Z$ is the same as the
cohomology of $S^2$ by factoring out the torus. The generator of
$H^2(Z)$ will be called $z_{11}$ with $z_{11}^2 = 0$.

All in all, this shows
\[ H^q(X_{\emptyset}) =
  \begin{cases}
    \Z & q = 0 \\
    \Z^7 & q = 1 \\
    \Z^{22} & q = 2 \\
    \Z^{42} & q = 3 \\
    \Z^{56} & q = 4 \\
    \Z^{56} & q = 5 \\
    \Z^{42} & q = 6 \\
    \Z^{22} & q = 7 \\
    \Z^{7} & q = 8 \\
    \Z & q = 9 \\
    0 & \text{otherwise}
  \end{cases} \] 
and the generators are
\[ H^*(X_{\emptyset}) =
\bigwedge\nolimits^*[a_{11},d_{11},g_{11},h_{11},i_{11},C_{11},D_{11}]
\otimes H^*(Z) \]

\subsubsection{A note on the cohomology of $Z$}

The space $Z$ defined above is a manifold, which we will prove by
applying the implicit function theorem. Consider the formula defining
$Z$,
\[ (x,y,z) \mapsto y-z-xyz-1 \]
Taking partial derivatives gives the three functions
\begin{align*}
  (x,y,z) &\mapsto yz \\
  (x,y,z) &\mapsto 1-xz \\
  (x,y,z) &\mapsto -1-xy
\end{align*}
If the first function is zero, then at least one of the two others
must be non-zero. By the implicit function theorem, the coordinate
that has non-zero derivative can be written as a smooth function of
the other coordinates. This shows that $Z$ is locally the graph of a
smooth function, so it is a smooth manifold.

Later on it will be necessary to have a description of $H^2(Z)$. This
is done by considering the subspace
\[ R = \set{(x,y,z) \in \R^3 \sdel x > 0, y < 0, xy < -1, z =
  \frac{y-1}{1+xy}>0 } \subset Z \]
This space is closed in $Z$ since any sequence in $R$ that converges
in $Z$ must also converge in $R$.
Consider a subbundle $NR \cong R\times D$ of the normal bundle of $R$
in the tangent bundle $TZ$, such that the exponential map is a
diffeomorphism on $NR$. This gives us a diagram
\[ \xymatrix{ H^2(NR,NR_0) & H^2(Z,Z-R) \ar[l]_\cong \ar[r]^>>>>>{j^*} &
  H^2(Z) } \]
where $NR_0$ is $NR$ without the zero section, $NR_0 \cong
R\times(D-0)$. The isomorphism is given by excision, since $Z = NR
\cup (Z-R)$. The space $R$ is homeomorphic to $\R^2$, so we have 
\[ H^2(Z,Z-R) \cong H^2(NR,NR_0) \cong H^2(R\times (D,D-0)) \cong
\Z \] 
From the diagram, $j^*(1)$ gives us an element of $H^2(Z)$. To see
that this is a generator, consider the torus 
\[ \widehat T = \set{(x,y,z)\in Z \mid \size{x} = \size{y} = 2 }
\subset Z \]
This intersects $R$ in exactly one point, namely $p=(2,-2,1)$. By
taking the intersection with $\widehat T$ we can extend the previous
diagram,
\[ \xymatrix{ H^2(NR,NR_0) \ar[d]^{i^*} & H^2(Z,Z-R) \ar[d]^{i^*}
  \ar[l]_<<<<<\cong \ar[r]^>>>>>{j^*} & H^2(Z) \ar[d]^{i^*} \\
  H^2(p\times (D,D-0)) & H^2(\widehat T,\widehat T-p)
  \ar[r]^>>>>>{\cong} \ar[l]_>>>>>\cong & H^2(\widehat T) 
} \]
The left $i^*$ is an isomorphism since the cohomology in degree two is
given by the pair $(D,D-0)$, which does not change under the
inclusion. This shows that the inclusion $\widehat T \to Z$ is an
isomorphism in cohomology, so $j^*(1)$ is a generator of
$H^2(Z)$. Taking the same diagram with cohomology replaced with
homology reverses all arrows but is otherwise exactly the same. This
shows that to calculate the homology class of a cycle $\sigma$ in
$H_2(Z)$, we can count the number of points of intersection between
$\sigma$ and $R$ with orientation, since each of these will contribute
either $1$ or $-1$, depending on orientation, to the homology class in
$(Z,Z-R)$.

There is a smiliar argument with the space $Y$ from before and the
subspace $R_Y = \set{(b,c) \in \R^2 \mid b > 0, c > 0, bc > 1}$. In
particular, the generator of $H_2(Y)$ can be taken as
$i_*\left([\widehat T]\right)$, where $[\widehat
  T]$ is a generator of $H_2(\widehat T)$ for the torus $\widehat T =
  \set{(b,c)\mid \size{b}=\size{c} = 2}$.

\subsection{The first page}

With all of the above, we can draw the $E_1$-page of the spectral
sequence. Remember that this is
\[ E_1^{p,q} \cong
\begin{cases}
  H^{q}(X_\emptyset) & p=0 \\
  H^{q+1}(X_{\set{1}} \times (D,D-0)) \oplus H^{q+1}(X_{\set{2}} \times
  (D,D-0)) & p=1 \\
  H^{q+2}(X_{\set{1,2}}\times (D_2,D_2-0)) & p = 2 \\
  0 & \text{otherwise}
\end{cases} \]

By inserting the groups that were calculated above into this, we get
Figure \ref{e1} with the non-zero differentials moving directly to
the right. To get the $E_2$-page of the spectral sequence, we need to
compute the homology of the differentials.

\begin{figure}[ht]
  \[ \xy
  (5,4); (5,-80); **\dir{-}; (5,4); *\dir{>};
  (5,-80); (42,-80); **\dir{-}; *\dir{<};
  (0,-82)*{q}; (5,-84)*{p};
  (1,-84); (4,-81); **\dir{-};
  (0,0)\xymatrix@R=1em@C=1.7em{
    9 & \Z & 0 & 0 \\
    8 & \Z^7 \ar[r]& \Z^2 & 0 \\
    7 & \Z^{22} \ar[r]& \Z^{12} \ar[r]& \Z \\
    6 & \Z^{42} \ar[r]& \Z^{32} \ar[r]& \Z^{4} \\
    5 & \Z^{56} \ar[r]& \Z^{50} \ar[r]& \Z^{7} \\
    4 & \Z^{56} \ar[r]& \Z^{50} \ar[r]& \Z^{7} \\
    3 & \Z^{42} \ar[r]& \Z^{32} \ar[r]& \Z^{4} \\
    2 & \Z^{22} \ar[r]& \Z^{12} \ar[r]& \Z \\
    1 & \Z^{7} \ar[r]& \Z^{2} & 0 \\
    0 & \Z & 0 & 0 \\
    & 0 & 1 & 2
  } \endxy \]
  \caption{The $E_1$-page with differentials.}
  \label{e1}
\end{figure}

\section{Differentials on $E_1$}

By definition, the differentials are the boundary maps of the long
exact sequence of the triple $(F_{p+1},F_p,F_{p-1})$,
\[ d_1 : E_1^{p,q} = H^{p+q}(F_p,F_{p-1}) \to H^{p+q+1}(F_{p+1},F_p) =
E_1^{p+1,q} \]
We will consider the two cases $p=0$ and $p=1$ separately.

\subsection{First column}

For $p = 0$, we are looking at the boundary map of the pair
$(F_1,F_0)$. Consider the inclusion of pairs,
\[ i : X_{\set{1}}\times (D,D-0) \sqcup X_{\set{2}}\times (D,D-0) \to
(F_1,F_0) \]
This is an isomorphism in cohomology and gives a map between the long
exact sequences that makes the following diagram
commute:
\[ \xymatrix{
H^{p+q}(F_0) \ar[r]^>>>>>>>>>>>>>>>>>>>>{d_1} \ar[d]^{i^*} & H^{p+q+1}(F_1,F_0)
\ar[d]^{i^*}_{\cong} \\
H^{p+q}((X_{\set{1}}\sqcup X_{\set{2}})\times (D-0))
\ar[r]^>>>>>{\delta} & H^{p+q+1}((X_{\set{1}}\sqcup X_{\set{2}})\times (D,D-0))
} \]
The bottom map is the boundary map from the long exact sequence of the
pair $(D,D-0)$. This map is well-known and will be treated at the end
of this section, so the only missing part is to calculate the
inclusion $i^*$. Since we are working with a disjoint union, we only
need to calculate the inclusion on each factor.

\subsubsection{The inclusion $X_{\set{1}}\times (D-0) 
  \hookrightarrow X_\emptyset$}

We can write the inclusion out in coordinates as
\begin{align*}
  (\C^*)^5\times Y \times (D-0) \longrightarrow
  (&\C^*)^7\times Z \\
  (a,d,f,h,i,b,c,g) \mapsto \bigg(&g,h,i,\frac{a}{g},-f,d,\\
  & C =
  c-\frac{a}{g}-\frac{cd}{f}-bc^2+\frac{abc}{g}, \\
  & \frac{bfC}{d},\frac{c}{C}-\frac{a}{g C},-\frac{cd}{fC}\bigg) 
\end{align*}

On the generators, this is given by
\begin{align*}
  i^* : H^*(X_{\emptyset}) &\to H^*(X_{\set{1}}\times (D-0)) \\
  a_{11} & \mapsto a_{01} - g_{01} \\
  d_{11} &\mapsto f_{01} \\
  g_{11} &\mapsto g_{01} \\
  h_{11} &\mapsto h_{01} \\
  i_{11} &\mapsto i_{01} \\
  D_{11} &\mapsto d_{01} \\
  C_{11} &\mapsto a_{01} - g_{01} + y_{01,1} \\
  z_{11} &\mapsto y_{01,2} + (a_{01}-d_{01}+f_{01}-g_{01})y_{01,1}
\end{align*}
where $g_{01}$ is the generator of $H^1(D-0)$.

In degree 1, this is calculated by dualizing to homology, writing out
the map on generators and working
out what they map to. For $z_{11}$ it is slightly
harder. To calculate this, we use the submanifold $R \subset Z$
defined previously as
\[ R = \set{(x,y,z)\in \R^3 \delim \Big\vert\delim x > 0, y < 0, z =
  \frac{y-1}{1+xy}>0} \]
by choosing a generator $\alpha = [\sigma]$ of $H_2(X_{\set{1}}\times
(D-0))$ and calculating the intersection of the image of $f = \pr_Z
\circ i\circ \sigma$ and $R$. Note that if the coordinates of $Y$ are
fixed at $(0,1)$, then the image of
\[ f(a,d,f,h,i,0,1,g) =
\left(0,\frac{1}{C}-\frac{a}{gC},-\frac{d}{fC} \right)\]
never intersects $R$ since $x = 0$. Likewise, the coordinates $i$ and
$h$ do not appear in the projection to $Z$ and can not influence the
result. Hence we only have to check the duals of the following
generators in degree 2:
\[ a_{01} y_{01,1}, d_{01} y_{01,1}, f_{01} y_{01,1}, y_{01,1} g_{01},
y_{01,2} \]
These will be handled individually. All non-relevant coordinates will
be set to 1, except $g$ which will be denoted $\e$ to make it clear
that this coordinate is small. If $y_{01,1}$ is involved, we choose the
generator
\[ \sigma(\lambda) = (1-\lambda,1) \qquad \lambda\in S^1 \]
as this gives
\[ y_{01,1} = [t \mapsto 1-(1-\lambda)\cdot 1] = [t\mapsto \lambda] \]

\paragraph{$\mathbf{a_{01} y_{01,1}}$:}
The relevant map is
\[ (a,\lambda) \mapsto \left( (1-\l)C, \frac{1}{C}-\frac{a}{\e
    C},\frac{1}{C}\right)\]
where $C = -1-\frac{a\l}{\e}$. If this intersects $R$ then $C$ has
to be real and positive by looking at the last coordinate, which
forces $\l$ real and negative by the first coordinate and $a$ real and
positive by the second. Hence there is at most one intersection, at
$(1,-1)$:
\[ f(1,-1) = \left( \frac{2-2\e}{\e},-1,\frac{\e}{1-\e}\right) \]
By inspection, this satisfies the conditions and hence belongs to
$R$, so we conclude that the image intersects once. This gives the
following formula:
\[ \langle f^*(z_{11}), a_{01}y_{01,1} \rangle = \langle
z_{11},f_*(a_{01}y_{01,1})\rangle = \pm 1 \]
To check the sign, we need to check if this map preserves the
orientation of the torus in $R$. This is done by differentiating the
map to get a map of tangent spaces
\[ Tf : TT^2 \to TR \]
and checking that the determinant of this map is
positive. This has been done in Sage.

\paragraph{$\mathbf{d_{01} y_{01,1}}$:}
In this case the map becomes
\[ (d,\l) \mapsto \left( \frac{(1-\l)C}{d},\frac{\e-1}{\e C},
  \frac{d}{C} \right) \]
and by a similar analysis, this intersects in exactly the
point $f(1,-1)$, which is the same as before. However, this time the
orientation is reversed.

\paragraph{$\mathbf{f_{01} y_{01,1}}$:}
This has the same point of intersection and preserves the orientation.

\paragraph{$\mathbf{y_{01,1}g_{01}}$:}
This also intersects and preserves the orientation.

\paragraph{$\mathbf{y_{01,2}}$:}
This is calculated by the intersection points of the fundamental class
of the torus in $Y$, which was mentioned earlier:
\[ \widehat T = \set{(b,c) \in \C^2 \mid \size{b}=\size{c} = 2}
\subset Y \]
The map becomes
\[ (b,c) \mapsto \left( bC, \frac{c}{C}-\frac{1}{\e C},\frac{c}{C}
\right) \]
with $C = -bc^2-\frac{1-bc}{\e}$. We will write $x$ for the first
coordinate, $y$ for the second and $z$ for the third. An analysis as
before gives $bc = xz > 0$ and $c = \frac{1}{\e(z-y)} \in
\R$. So both $b$ and $c$ must be real and $c$ should be positive. Then
$xy = bc - b\e^{-1}<0$ shows that $b$ must be positive since $bc$ is
positive. The only possible point of intersection is
\[ f(2,2) = \left( \frac{6-16\e}{\e},
  \frac{2\e-1}{3-8\e},\frac{2\e}{3-8\e} \right) \]
Since $\e$ is very small, this intersects $R$ and the orientation is
preserved.

\subsubsection{The inclusion $X_{\set{2}}\times (D-0) 
  \hookrightarrow X_\emptyset$}

The inclusion in coordinates is
\begin{align*}
  (\C^*)^5\times Y \times (D-0) \longrightarrow (&\C^*)^7\times Z \\
  (a,d,f,g,i,b,c,h) \mapsto 
  \bigg(g&,h,i,a,-\frac{d}{h},\\
  & D=\frac{d}{h}(1-bc), \\
  & C= -af-\frac{a^2bh}{d}+\frac{a^2bhf}{d}+\frac{a^2b^2ch}{d}, \\
  & \frac{bC}{D}, \frac{cd}{hC}-\frac{a}{C},
  \frac{cD}{C}-\frac{ahD}{dC}+\frac{afhD}{dC}+
  \frac{abchD}{dC}\bigg)
\end{align*}

We calculate as before and get the following on generators:
\begin{align*}
  i^* : H^*(X_{\emptyset}) &\to H^*(X_{\set{2}}\times (D-0)) \\
  a_{11} & \mapsto a_{10} \\
  d_{11} &\mapsto d_{10} - h_{10} \\
  g_{11} &\mapsto g_{10} \\
  h_{11} &\mapsto h_{10} \\
  i_{11} &\mapsto i_{10} \\
  D_{11} &\mapsto d_{10} + y_{10,1} - h_{10} \\
  C_{11} &\mapsto f_{10} + a_{10} \\
  z_{11} &\mapsto y_{10,2} + (a_{10}-d_{10}+f_{10}+h_{10})y_{10,1}
\end{align*}
where $h_{10}$ is the generator of $H^1(D-0)$.

The calculations have been left out, but are very similar to the
previous ones.

\subsection{Second column}

The differential on the second column is the boundary map
\[ \delta : H^q((X_{\set{1}}\sqcup X_{\set{2}})\times (D,D-0)) \cong
H^q(F_1,F_0) \to H^{q+1}(F_2,F_1) \] 

To calculate this map, consider the inclusion of triples
\[ X_{\set{1,2}}\times (D_2,D_2-0,(D-0)\times (D-0)) \hookrightarrow
(F_2,F_1,F_0) \]
given by inserting a small disc, sphere or torus on the zero-entries
of $X_{\set{1,2}}$. There are homeomorphisms of pairs,
\begin{align*}
  \big(D_2-0,(D-0)\times(D-0)\big) \cong \big(D\times (D-0) \sqcup
  (D-0)\times D, (D-0)\times(D-0)\big) \\
  \big(D_2,D_2-0\big) \cong \big(D\times D,D\times
  (D-0)\cup (D-0)\times D\big) =  (D,D-0)\times (D,D-0)
\end{align*}
These gives a commuting diagram,
\[ \xymatrix{ H^q(F_1,F_0) \ar[r] \ar[d] & H^{q+1}(F_2,F_1)
  \ar[d]^\cong \\ 
  H^q(X_{\set{1,2}}\times (D_2-0,(D-0)\times (D-0))) \ar[r]
%  \ar@{=}[d] 
&
  H^{q+1}(X_{\set{1,2}} \times (D_2,D_2-0)) %\ar@{=}[d] 
%\\
%  H^q(X_{\set{1,2}}\times ((D,D-0)\times (D-0) \sqcup (D-0)\times
%  (D,D-0))) \ar[r] & H^{q+1}(X_{\set{1,2}}\times (D,D-0)\times (D,D-0))
} \]
The marked isomorphism follows from excision and the left map is
induced by an inclusion. If we apply the
homeomorphism of pairs to the bottom of the diagram, the boundary map
at the bottom becomes the sum of the two boundary maps,
\[ \xymatrix@R=0.1em{ H^q(X_{\set{1,2}}\times (D,D-0)\times (D-0))
  \ar[]!<15ex,0ex>;[rd]!<-20ex,0ex>^{\delta} \\
  & H^{q+1}(X_{\set{1,2}}\times (D,D-0)\times (D,D-0)) \\
  H^q(X_{\set{1,2}}\times (D-0)\times (D,D-0))
  \ar[]!<15ex,0ex>;[ur]!<-20ex,0ex>_\delta } \]
The boundary maps are given by the long exact sequence of the pair
$(D,D-0)$, so we only need to consider the inclusion on the left. By
choosing the discs sufficiently small and using the identification of
$H^q(F_1,F_0)$ defined in the first section, 
we see that the inclusion maps the space $X_{\set{1,2}}\times
(D,D-0)\times (D-0)$ to $X_{\set{1}}\times(D,D-0)$ and
$X_{\set{1,2}}\times (D-0)\times (D,D-0)$ to $X_{\set{2}}\times
(D,D-0)$. To get the differential, we must calculate the sum of these
two maps:
\begin{align*}
  i_1^* : H^q(X_{\set{1}}\times (D,D-0)) &\to H^q(X_{\set{1,2}}\times
  (D,D-0)\times (D-0)) \\
  i_2^* : H^q(X_{\set{2}}\times (D,D-0)) &\to H^q(X_{\set{1,2}}\times
  (D-0) \times (D,D-0))
\end{align*}

\subsubsection{The inclusion $X_{\set{1,2}}\times (D,D-0) 
\times (D-0) \to X_{\set{1}} \times (D,D-0)$}   

In coordinates, the map becomes
\begin{align*}
  (\C^*)^3 \times Y \times (D,D-0)\times (D-0) &\to (\C^*)^5 \times Y
  \times (D,D-0) \\
  (a,d,i,b,c,g,h) &\mapsto \left(a,i,h,\frac{d}{h},
  \frac{d}{h}(1-bc), \frac{bh}{d},\frac{dc}{h},g\right)
\end{align*}

Again, we calculate the map on cohomology by dualizing to homology
and working with the dual generators. To check what happens to the
elements of $H^q(Y)$, we count intersections with the subspace $R_Y =
\set{(b,c)\in \R^2 \mid b > 0, c > 0, bc > 1}$ of $Y$, exactly as for
$Z$ above.
\begin{align*}
  i_1^* : H^q(X_{\set{1}}\times (D,D-0)) &\to H^q(X_{\set{1,2}}\times
  (D,D-0)\times (D-0)) \\
  a_{01} &\mapsto a_{00} \\
  d_{01} &\mapsto d_{00}-h_{00} + y_{00,1} \\
  f_{01} &\mapsto d_{00} - h_{00} \\
  g_{01} &\mapsto g_{00} \\
  h_{01} &\mapsto h_{00} \\
  i_{01} &\mapsto i_{00} \\
  y_{01,1} &\mapsto y_{00,1} \\
  y_{01,2} &\mapsto y_{00,2} - d_{00}y_{00,1} + h_{00}y_{00,1}
\end{align*}
We denote the generator of $H^2(D,D-0)$ in
$H^*(X_{\set{1}}\times (D,D-0))$ by $g_{01}$ and for
$H^*(X_{\emptyset} \times (D,D-0) \times (D-0))$ we consider $g_{00}\in
H^2(D,D-0)$ and $h_{00} \in H^1(D-0)$. 

\subsubsection{The inclusion $X_{\set{1,2}}\times (D-0)
\times (D,D-0) \to X_{\set{2}}\times (D,D-0)$}
This is exactly as above for the map
\begin{align*}
  (\C^*)^3 \times Y \times (D-0)\times (D,D-0) &\to (\C^*)^5 \times
  Y\times (D,D-0) \\ 
  (a,d,i,b,c,g,h) &\mapsto
  \left(\frac{a}{g},d,i,g,1-bc,b,c,h\right)
\end{align*}
In cohomology, this becomes
\begin{align*}
  i_2^* : H^q(X_{\set{2}}\times (D,D-0)) &\to H^q(X_{\set{1,2}}\times
  (D-0) \times (D,D-0)) \\
  a_{10} &\mapsto a_{00} - g_{00} \\
  d_{10} &\mapsto d_{00} \\
  f_{10} &\mapsto y_{00,1}\\
  g_{10} &\mapsto g_{00} \\
  h_{10} &\mapsto h_{00} \\
  i_{10} & \mapsto i_{00} \\
  y_{10,1} &\mapsto y_{00,1} \\
  y_{10,2} &\mapsto y_{00,2}
\end{align*}
with $g_{01}\in H^2(D,D-0)$ inside $H^*(X_{\set{2}}\times (D,D-0))$ and
$g_{00}\in H^2(D,D-0)$, $h_{00}\in H^1(D-0)$ inside
$H^*(X_{\emptyset}\times(D-0)\times(D,D-0))$.

\subsubsection{Boundary maps}
\label{disk}
We have now reduced everything to calculating some well-known boundary
maps for the pair $(D,D-0)$. By writing out the long exact sequence,
the boundary map
\[ \delta : H^*(D-0) \to H^*(D,D-0) \]
is an isomorphism in degree 1,
\[ H^1(D-0) \cong H^2(D,D-0) \]
and zero in all other degrees. The boundary map for the pair $X\times
(D,D-0)$ is $\Id\otimes \delta$ in tensor-notation. So the boundary
map in the cases above picks out the ``new'' coordinate added to
$X\times (D-0)$ and throws away everything that does not depend on
it. To be more precise, the map is given on generators by:
\begin{align*}
  H^*(X)\otimes H^*(D-0)) &\to H^*(X)\otimes H^*(D,D-0)) \\
  x\otimes d &\mapsto
  \begin{cases}
    x \otimes \delta d & d\in H^1(D-0) \\
    0 & \text{otherwise}
  \end{cases}
\end{align*}

\section{Turning the page}

We can now use the above to calculate the $E_2$-page. We have
calculated the groups appearing on the first page and the
differential. To get the second page, we need to compute the homology
groups of $d_1$. This reduces to linear algebra over the integers, and
is done using Sage. The result can be found in Figure \ref{e2}.

\begin{figure}[ht]
  \[ \xy
  (5,4); (5,-80); **\dir{-}; (5,4); *\dir{>};
  (5,-80); (42,-80); **\dir{-}; *\dir{<};
  (0,-82)*{q}; (5,-84)*{p};
  (1,-84); (4,-81); **\dir{-};
  (0,0)\xymatrix@R=1em@C=1.7em{
    9 & \Z & 0 & 0 \\
    8 & \Z^5 & 0 & 0 \\
    7 & \Z^{11} & 0 & 0 \\
    6 & \Z^{15} & \Z & 0 \\
    5 & \Z^{16} & \Z^{3} & 0 \\
    4 & \Z^{16} & \Z^{3} & 0 \\
    3 & \Z^{15} & \Z & 0 \\
    2 & \Z^{11} & 0 & 0 \\
    1 & \Z^{5} & 0 & 0 \\
    0 & \Z & 0 & 0 \\
    & 0 & 1 & 2
  } \endxy \]
  \caption{The $E_2$-page.}
  \label{e2}
\end{figure}

The differentials from $E_2$ onwards are given as maps:
\[ d_r : E_r^{p,q} \to E_r^{p+r,q+1-r} \]
For $r\geq 2$, this increases $p$ by at least 2. This means that
the domain or codomain of $d_r$ will always be zero, so taking
homology does not change anything. Hence $E_2 = E_\infty$, the
spectral sequence collapses at the 2\textsuperscript{nd} term, and
there are no extension problems since all the groups are free and
abelian. This allows us to read off the cohomology of $X_{3,3}$. Since
we are also interested in factoring out the $T^3$ torus, we note this
as well.
\[ H^q(X_{3,3}) = E_2^{0,q} \oplus E_2^{1,q-1} = 
\begin{cases}
  \Z & q = 0 \\
  \Z^{5} & q = 1\\
  \Z^{11} & q = 2\\
  \Z^{15} & q = 3\\
  \Z^{17} & q = 4\\
  \Z^{19} & q = 5\\
  \Z^{18} & q = 6\\
  \Z^{12} & q = 7\\
  \Z^{5} & q = 8\\
  \Z & q = 9\\
  0 & \text{otherwise} 
\end{cases} \qquad
H^q(Y_{3,3}) =
\begin{cases}
  \Z & q = 0 \\
  \Z^{2} & q=1 \\
  \Z^{2} & q=2 \\
  \Z^{2} & q=3 \\
  \Z^{3} & q=4 \\
  \Z^{2} & q=5 \\
  \Z & q=6 \\
  0 & \text{otherwise}
\end{cases}
\]
We can also read off the ring structure from the spectral sequence, as
it is induced from the spaces we start with. Since we do not
need it, it will not be included here.
% Ogs� noget med \pi_1? Eg. surjektivitet eller lignende?


%%% Local Variables: 
%%% mode: latex
%%% TeX-master: "main"
%%% End: 
 % Bliver nok langt.... Stjæl fra noter

\chapter{Loop spaces}

We now return to the general case. We want to show that the space
$Y_{m,\infty}$ is equivalent to the space of loops in the flag
manifold $\GLB{m}$.

\begin{definition}
  The loop space $\Omega = \Omega(\GLB{m})$ is the space
  \[ \Omega = \set{\g : [0,1] \to \GLB{m} \mid \g(0) = \g(1) =
    [\Id]} \]
  with the compact-open topology. 
\end{definition}

\begin{remark}
  Following \cite[Chapter 17]{milnor}, the compact-open topology is
  also the topology given by the metric
  \[ d(\g,\g') = \max_{t\in[0,1]} \widetilde d (\g(t),\g'(t)) \]
  where $\widetilde d$ is a metric on $\GLB{m}$ induced from a
  Riemannian metric. The space $\Omega$ is homotopy equivalent to the
  space of piecewise smooth loops with a slightly modified metric, and
  we will generally choose the most convenient model for our
  arguments.
\end{remark}

\section{The map $Y_{m,n}\to\Omega$}

For $A = [a_1,\dots,a_n]\in \X$ we
can define a path $ \g_A$ in $\GL_m$ by starting with
$\g_A(0) = \Id$. Since the matrix $[e_2,\dots,e_m,a_1]$ is
invertible by the definition of $\X$, we know that $a_1$ is linearly
independent of $e_2,\dots,e_m$. Hence it can be written as
\[ a_1 = \sum_{i=1}^m c_i e_i \]
with $c_1 \neq 0$. Rewriting the expression leads to the formula
\[  \frac{1}{c_1} a_1= e_1 + \frac{1}{c_1} \left(\sum_{i=2}^m c_i
  e_i\right) \]
This gives us a path in $\GL_m$ from the identity to
$\left[\frac{1}{c_1} a_1,e_2,\dots,e_m\right]$:
\[ [0,1]\ni t \mapsto \left[ e_1 + t\left(\sum_{i=2}^m
    \frac{c_i}{c_1}e_i\right),e_2,\dots,e_m \right] \in \GL_m \]
This process can be continued, since we can write
\[ \frac{1}{d_2}a_2 = e_2 + \frac{1}{d_2}\left(\sum_{i=3}^m d_i
  e_i\right) + \frac{d_1}{d_2c_1}a_1 \]
and we get a path
\[ t\mapsto \left[\frac{1}{c_1}a_1,e_2+t\left(\sum_{i=3}^m
    \frac{d_i}{d_2}e_i +
    \frac{d_1}{d_2c_1}a_1\right),e_3,\dots,e_m\right] \]
that starts where the other path stopped. If we continue through the
vectors in $A$, we will eventually end at a point
\[ [\l_i e_i,\dots,\l_me_m,\l_1e_1,\dots,\l_{i-1}e_{i-1}]\in \GL_m \]
which is diagonal, except the columns may have been
permuted if $n$ is not divisible by $m$. Choosing a path in $\GL_m$
from $[e_i,\dots,e_{i-1}]$ to $\Id$, scaling the columns and appending
it to all the other paths then gives a path $ \g_A$ with
\[  \g_A(0) = \Id, \qquad \g_A(1) =
[\l_1e_1,\dots,\l_me_m] \]

If we consider this as a path in $\GLB{m}$ we
get a loop, since $\g_A(1)$ is a diagonal matrix. The
assignment
\[ f(A) = \g_A \in \Omega(\GLB{m}) \]
defines a continuous function from $\X$ to the loop space. Note that
this function descends to the quotient space $Y_{m,n}=\X/T^n$, since
elements of $\X$ that are equivalent under the $T^n$-action gives
paths in $\GL_m$ that are equivalent under the $B_m$-action.

The goal of this chapter is to prove the following theorem:

\begin{theorem}
  The map $f : Y_{m,n} \to \Omega$ induces a map
  \[ f : Y_{m,\infty} \to \Omega \]
  This is a (weak) homotopy equivalence.
\end{theorem}

Since the space $\Omega$ is relatively well-studied, this gives
information about $Y_{m,\infty}$. In particular, it allows us to
calculate the homology and cohomology of this space.

We want to show that the stabilization map
\begin{align*}
  s : Y_{m,n} &\to Y_{m,n+1} \\
  A = (a_1,\dots,a_n) &\mapsto \left(La_1,\dots,La_n,
    \begin{pmatrix}
      1 \\
      \vdots\\
      1
    \end{pmatrix} 
\right) = (LA,Le_1)
\end{align*}
does not affect the map $f$, or more precisely that $f$ induces a
well-defined map on $Y_{m,\infty}$. To do this we will work with the
subsequence $Y_{m,nm}$ of the sequence $Y_{m,n}$, as the two limits
are isomorphic.

We will show that the diagram
\[ \xymatrix{ Y_{m,nm} \ar[rr]^{s^m}\ar[dr]_f & & Y_{m,(n+1)m} \ar[dl]^f \\
  & \Omega & } \]
commutes up to homotopy. This is done by showing that the map $s^m$ is
homotopic to the map $\widetilde s$ that inserts the identity on the
last $m$ columns of $A \in Y_{m,nm}$:
\[ \widetilde s (a_1,\dots,a_{nm}) = (a_1,\dots,a_{nm},e_1,\dots,e_m) \]
If we connect $\Id$ and $L^{-m}$ by a path $\g$ in the space of
invertible lower triangular matrices, we get a homotopy
\[ s_t(A) = (\g(t)L^mA,
\g(t)L^me_1,\g(t)L^{m-1}e_1,\dots,\g(t)Le_1) \]
which goes from $s_0 = s^m$ to the map
\[ s_1 : A \mapsto (A,e_1,L^{-1}e_1,\dots,L^{-m+1}e_1) \]
The matrix $T = [e_1,L^{-1}e_1,\dots,L^{-m+1}e_1]$ is upper
triangular with entries given by the formula:
\[ T_{ij} = (L^{-j+1}e_1)_i = (-1)^{i-1}\binom{j-1}{i-1} \quad
\text{when } i \leq j \]
This is proven by considering negative powers of $L$ and showing the
identity:
\[ (L^{-j+1})_{ik} = (-1)^{i-k}\binom{j-1}{i-k} \]
But since $T$ is an upper triangular matrix, it can be moved to the
identity without affecting the linear independence.
Hence the maps $s^m$ and $\widetilde s$ are homotopic. The composition
$f\circ \widetilde s$ is the same as $f$, except that the loop
$f\circ\widetilde s$ is stationary at the identity for a while before
it terminates. Putting it all together shows the following:
\begin{lemma}
  The diagram
  \[ \xymatrix{ Y_{m,nm} \ar[rr]^{s^m}\ar[dr]_{f} & & \ar[dl]^f
    Y_{m,(n+1)m} \\
    & \Omega & } \]
  commutes up to homotopy and $f$ induces a map:
  \[ f : Y_{m,\infty} \to \Omega \]
\end{lemma}

\section{Homotopy equivalence}

To show that the map $f : Y_{m,\infty} \to \Omega$
gives an isomorphism of homotopy groups, we will again follow
\cite{milnor}. For each element $U$ of $\SUT{m}$, we can find a 
geodesically convex open ball centered on $U$. Since $\SUT{m}$
is compact, there is some $\e > 0$ such that any ball of radius $\e$
or less is geodesically convex. By the same argument, we can also
choose $\e$ so small that for any two points $A= [a_1,\dots, a_m]$ and
$B=[b_1,\dots,b_m]$, with the distance from $A$ to $B$ less
than $\e$, any $m$ subsequent vectors in the sequence
$(a_1,\dots,a_m,b_1,\dots,b_m)$ are linearly independent.
With this, we get an increasing sequence of open sets in $\Omega$,
defined by requiring the loops to be piecewise contained in a ball of
radius $\e$:
\[ \Omega_k = \set{\g \in \Omega \delim \Big\vert\delim
  \g_{\left|\left[ \frac{j-1}{2^k},\frac{j}{2^k} \right]\right.} \text{ is
  contained in an $\e$-ball $B_\e$} } \]
By the choice of $\e$ made above, we can define a map
\begin{align*}
  \varphi_k : \Omega_k &\to Y_{m,(2^k-1)m} \\
  \g &\mapsto \left[ \g\left(\frac{1}{2^k}\right),
    \g\left(\frac{2}{2^k}\right), \dots,
    \g\left(\frac{2^k-1}{2^k}\right) \right]
\end{align*}
Note that $\g(t)$ is an element of $\SUT{m}$, so the notation
$[\g(t),\dots]$ is shorthand for $[\g(t)e_1,\dots,\g(t)e_m,\dots]$.
The map is continuous since two paths being close in $\Omega_k$ means
that the maximum distance between them is small. But this translates
to the vectors of the image being close together. The maps $\varphi_k$
and $\varphi_{k+1}$ are related in the following fashion:

\begin{lemma}
  \label{lem:com}
  The following diagram commutes up to homotopy:
  \[ \xymatrix{\Omega_k \ar@{^(->}[r] \ar[d]^{\varphi_k} & \Omega_{k+1}
  \ar[d]^{\varphi_{k+1}} \\
  Y_{m,(2^k-1)m} \ar[r]^{s} & Y_{m,(2^{k+1}-1)m} } \]
\end{lemma}
\begin{proof}
  The proof is by direct computation.
  \begin{align*}
    \varphi_{k+1} &= \left( \g \mapsto \left[
        \g\left(\frac{1}{2^{k+1}}\right),
        \g\left(\frac{2}{2^{k+1}}\right), \dots,
        \g\left(\frac{2^{k+1}-1}{2^{k+1}}\right) \right]\right) \\
    &\simeq \left(\g \mapsto \left[ \g\left(\frac{2}{2^{k+1}}\right),
      \g\left(\frac{2}{2^{k+1}}\right),
      \g\left(\frac{4}{2^{k+1}}\right), \dots,
      \g\left(\frac{2^{k+1}}{2^{k+1}}\right) \right]\right) \\
    &\simeq \left(\g\mapsto \left[ \g\left(\frac{1}{2^{k}}\right),
      \g\left(\frac{2}{2^{k}}\right), \dots,
      \g\left(\frac{2^k-1}{2^k}\right), \Id,\dots,\Id \right] \right)\\
    &= s\circ \varphi_k
  \end{align*}
  The first homotopy is given by sliding along the minimal geodesic
  between the points and gives an allowed path by our choice of
  $\e$. The second homotopy is somewhat similar. If $(a_1,\dots,a_m)$
  and $(b_1,\dots,b_m)$ are two elements of $\SUT{m}$ and we have an
  element of $Y_{m,n}$ of the form
  \[ \left[ \dots, a_1,\dots,a_m,a_1,\dots,a_m,b_1,\dots,b_m,\dots
  \right]\]
  then we can change the middle columns one by one, as we did in
  the construction of $f$, and get a path in $Y_{m,n}$ that changes
  the middle from $[a_1,\dots,a_m]$ to $[b_1,\dots,b_m]$:
  \begin{align*}
    \left[a_1,\dots,a_m,a_1,a_2,\dots,a_m,b_1,\dots,b_m \right] 
    &\leadsto
    \left[a_1,\dots,a_m,b_1,a_2,\dots,a_m,b_1,\dots,b_m\right] \\ 
    &\leadsto
    \left[a_1,\dots,a_m,b_1,b_2,\dots,a_m,b_1,\dots,b_m\right] \\  
    &\leadsto
    \left[a_1,\dots,a_m,b_1,b_2,\dots,b_m,b_1,\dots,b_m\right]  
  \end{align*}
  This allows us to take all the repetitions
  $\left[\dots,\g(t),\g(t),\dots\right]$ and move them to the end as
  diagonal matrices, which results in the stabilization map $s$.
\end{proof}

The above lemma shows that the maps $\varphi_k$ fit together to define
a map from $\Omega$ to $Y_{m,\infty}$.
The idea is to use the maps $\varphi_k$ as homotopy inverses of the
loop space map $f$. We are now ready to prove the result.

\begin{theorem}
  The map $f : Y_{m,\infty} \to \Omega$ is a weak homotopy
  equivalence, i.e. the map
  \[f_* : \pi_i(Y_{m,\infty}) \to \pi_i(\Omega)\]
  is an isomorphism for all $i$.
\end{theorem}

\begin{proof}
  We will start with surjectivity. Let
  \[ g : W \to \Omega \]
  be a finite cell complex in $\Omega$. We want to show that we can
  construct a map $\widetilde g$ from $W$ to $Y_{m,\infty}$ such that
  the composition $f\circ \widetilde g$ is homotopic to $g$.
  
  The image of $g$ is compact and hence contained in $\Omega_k$ for
  some $k$, since $\set{\Omega_k}$ is an open cover of
  $\Omega$. Consider the map 
  \[ \widetilde g = \varphi_k\circ g : W \to Y_{m,(2^k-1)m} \]
  Evaluating at a point $w\in W$ gives
  \[ \widetilde g(w) = \varphi_k(\g_w) = \left[
    \g_w\left(\frac{1}{2^k}\right),\dots,
    \g_w\left(\frac{2^k-1}{2^k}\right) \right] \]
  Evaluating $f$ on this gives a path connecting the points
  $\g_w\left(\frac{j}{2^k}\right)$. But by our choice of $\e$, any
  such path is homotopic to the path connecting these points with
  minimal geodesics, so $f\circ\widetilde g \simeq g$.

  To show injectivity, let
  \[ g : W \to Y_{m,\infty}\]
  be a finite cell complex with $f\circ g$ null-homotopic. We want to
  show that $g$ is null-homotopic in $Y_{m,\infty}$. In the
  direct limit, a compact set is contained in the image of $Y_{m,n}$
  for some $n$ (see \cite[Proposition~A.1]{hatcher} or
  \cite[Chapter~9.4]{may}), so $g$ is represented by a map:
  \[ g : W \to Y_{m,n} \]
  We will work with this representative. The map $f\circ
  g$ is null-homotopic in $\Omega$ and the image is compact, so it is
  contained in $\Omega_k$ for some $k$. By applying $\varphi_k$, we
  get a function
  \[ \varphi_k\circ f\circ g : W \to Y_{m,(2^k-1)m} \]
  Since $f\circ g$ is null-homotopic, this function is
  null-homotopic. But by the same argument that was used in
  Lemma \ref{lem:com}, the map
  \[ \varphi_k(f(g(w))) = \left[
    \g_w\left(\frac{1}{2^k}\right),\dots,\g_w\left(\frac{2^k-1}{2^k}
    \right)\right] \] 
  is homotopic to the stabilization map
  \[ s(g(w)) = \left[g(w),\Id,\dots,\Id\right] \]
  with the homotopy given by sliding along the path $\g_w=f(g(w))$.
  So the map $g$ can be stabilized to a null-homotopic map, so it must
  be null-homotopic in $Y_{m,\infty}$. This shows that $f_*$ is
  injective on homotopy groups.
\end{proof}

Since $f$ is a weak homotopy equivalence, it follows that $f_*$ is
an isomorphism on homology and cohomology groups. The Pontrjagin
homology ring of the loop space, with product given by loop
concatenation, has been calculated in \cite[Theorem~4.1]{grbic} as
\begin{align*}
  H_*(\Omega(\SUT{m})) \cong T(x_1,\dots,x_{m-1}) \otimes
  \Z[y_1,\dots,y_{m-1}]
\end{align*}
where $T(x_1,\dots,x_{m-1})$ is the tensor algebra generated by
$x_1,\dots,x_{m-1}$. The product has the relation $x_k^2 =
x_px_q+x_qx_p = 2y_1$ for $1\leq k,p,q\leq m-1$ and $p\neq q$. The
degree of $x_i$ is 1 and the degree of $y_i$
is $2i$. By the above result, this is also the homology of
$Y_{m,\infty}$. The groups are finitely generated and torsion free, so
by the Universal Coefficient Theorem the homology and cohomology
groups are isomorphic as abelian groups.


\section{Future work}

In the above, we showed that the space 
\[ Y_{m,\infty} = \varinjlim_n X_{m,n}/T^n \]
is equivalent to the loop space $\Omega(\SUT{m})$ and used this to
calculate the homology and cohomology groups of the limit. The
current aim of the project is to work with the spaces $Y_{m,n}$
directly and see what we can prove about these, in particular by using
the stabilization map
\[ s : Y_{m,n} \to Y_{m,n+1} \]
The conjecture at the moment is the following:
\begin{conjecture}
  The induced maps on cohomology
  \[ s^* : H^*(Y_{m,n+1}) \to H^*(Y_{m,n}) \]
  are always surjective and the induced maps on homology
  \[ s_* : H_*(Y_{m,n}) \to H_*(Y_{m,n+1}) \]
  are always injective.
\end{conjecture}
The idea for a proof would be to use a spectral sequence argument like
the one applied to $X_{3,3}$ by defining a filtration that is easier
to work with and then proceeding by induction, using that the maps
$s^*$ and $s_*$ induce maps of the spectral sequences.

Another thing to consider is modifying the space $\X$. One possible
change could be to replace $\X$ by a union:
\[ X'_{m,n} = \bigcup_{A\in\GL_m} X_{m,n}(A,A) \]
This would allow all sequences of vectors that start and end at the
same matrix, and one would expect the space $Y'_{m,\infty}$ to be
equivalent to the free loop space, consisting of all
maps from the circle $S^1$ to $\SUT{m}$.

There is also the option of replacing the complex numbers $\C$ with
some other space, like the reals or quaternions. The fact that $\R^*$
is not connected would most likely make the space $\X^{\R}$ quite
different, but possibly the space $\X^{\mathbb{H}}$ would behave in a
similar manner to $\X$, with some shifts in dimension.


%%% Local Variables: 
%%% mode: latex
%%% TeX-master: "main"
%%% End: 
 % Det jeg laver nu

\backmatter
%\nocite{*}
\printbibliography{}

\end{document}
